\uuid{0MDS}
\titre{Limite en un point}
\auteur{bodin ; hanani ; mzari}
\organisation{exo7}
\licence{CC BY-NC-SA 4.0}
\datecreate{2020-01-20}
\niveau{L1}
\module{Analyse}
\chapitre{Continuité, limite et étude de fonctions réelles}
\sousChapitre{Limite de fonctions}
\difficulte{3}
\theme{Limites des fonctions réelles, limite en un point}
\exo7qcmid{lille-1.limites.029}

\contenu{
\begin{question}
\qclassification{123.03}
\qtags{motcle=limite en un point}
Soit  $f$ la fonction définie sur $]\frac{3}{2}, +\infty[ \setminus \{2\}$  par : $f(x)=\left\{\begin{array}{cc}a\frac{\sqrt{x-1}-1}{x-2},& \mbox{si} \, x<2  \\ \frac{\sqrt{2x-3}-b}{x-2},&  \mbox{si} \, x >2  \end{array}\right.$. $f$ admet une limite  finie quand $x$ tend vers $2$ si et seulement si :
\begin{answers}

    \good{$a=2$ et $b=1$}

    \bad{ $a>0$ et $b >0$ }

    \bad{ $a=2$ et  $b >0$ }

    \bad{$a=0$ et $b=1$}

\end{answers}
\begin{explanations}
Si $b\neq 1$, $f$ admet une limite infinie quand $x$ tend vers $2^+$. On suppose que $b=1$ et on multiplie $f$ par l'expression conjuguée selon les  cas. On obtient :
$f(x)=\left\{\begin{array}{cc} \frac{a}{\sqrt{x-1}+1}& \mbox{si} \, x<2  \\ \frac{2}{\sqrt{2x-3}+1}&  \mbox{si} \, x >2  \end{array}\right.$. On déduit que $f$ admet une limite  finie quand $x$ tend vers $2$ si et seulement si $a=2$.
\end{explanations}
\end{question}
}
